%------------------------
% Resume Template
% Author : Harsh
% 
%------------------------

\documentclass[a4paper,20pt]{article}

\usepackage{latexsym}
\usepackage[empty]{fullpage}
\usepackage{titlesec}
\usepackage{marvosym}
\usepackage[usenames,dvipsnames]{color}
\usepackage{verbatim}
\usepackage{enumitem}
\usepackage[pdftex]{hyperref}
\usepackage{fancyhdr}
\usepackage{graphicx}
\usepackage{datetime2}
\usepackage{tikz}
\usepackage{blindtext}
\usepackage[document]{ragged2e}
\usepackage{xcolor}
\usepackage{hhline}
\usepackage{hyperref}
\usepackage{xcolor}

\hypersetup{
    colorlinks=true,
    linkcolor=blue,
    urlcolor=blue,
    citecolor=blue
}
\usepackage{tikz}

% Helper command to draw skill circles
\newcommand{\skilllevel}[1]{
  \begin{tikzpicture}[baseline=-0.5ex]
    \foreach \i in {1,...,5} {
      \ifnum\i>#1
        \draw[fill=white] (\i*0.25,0) circle (0.08);
      \else
        \draw[fill=black] (\i*0.25,0) circle (0.08);
      \fi
    }
  \end{tikzpicture}
}


\pagestyle{fancy}
\fancyhf{} % clear all header and footer fields
\fancyfoot{}
\renewcommand{\headrulewidth}{0pt}
\renewcommand{\footrulewidth}{0pt}

% Adjust margins
\addtolength{\oddsidemargin}{-0.530in}
\addtolength{\evensidemargin}{-0.375in}
\addtolength{\textwidth}{1in}
\addtolength{\topmargin}{-.45in}
\addtolength{\textheight}{1in}

\urlstyle{rm}

\raggedbottom
\raggedright
\setlength{\tabcolsep}{0in}

% Sections formatting
\titleformat{\section}{
  \vspace{-10pt}\scshape\raggedright\large
}{}{0em}{}[\color{blue}\titlerule \vspace{-6pt}]

%-------------------------
% Custom commands
\newcommand{\resumeItem}[2]{
  \item\small{
    \textbf{#1}{: #2 \vspace{-2pt}}
  }
}

\newcommand{\resumeItemWithoutTitle}[1]{
  \item\small{
    {\vspace{-2pt}}
  }
}

\newcommand{\resumeSubheading}[4]{
  \vspace{-1pt}\item
    \begin{tabular*}{0.97\textwidth}{l@{\extracolsep{\fill}}r}
      \textbf{#1} & #2 \\
      \textit{#3} & \textit{#4} \\
    \end{tabular*}\vspace{-5pt}
}

% Add this to your preamble
\usepackage{array}
\newenvironment{awlist}{% 
    \begin{tabular}{>{\bfseries}l @{\hspace{1em}} p{13cm}}
}{%
    \end{tabular}
}


\newcommand{\resumeSubItem}[2]{\resumeItem{#1}{#2}\vspace{-3pt}}

\renewcommand{\labelitemii}{$\circ$}

\newcommand{\resumeSubHeadingListStart}{\begin{itemize}[leftmargin=*]}
\newcommand{\resumeSubHeadingListEnd}{\end{itemize}}
\newcommand{\resumeItemListStart}{\begin{itemize}}
\newcommand{\resumeItemListEnd}{\end{itemize}\vspace{-5pt}}
%\newcommand{\resumeSubHeadingListStart}{\begin{itemize}[label={}, leftmargin=0pt]}
%\newcommand{\resumeSubHeadingListEnd}{\end{itemize}}

\usepackage{fontawesome5}
\usepackage{hyperref}
\hypersetup{colorlinks=true, urlcolor=blue}

%\usepackage{tgheros}
%\renewcommand{\familydefault}{\sfdefault}

\usepackage{charter}
%\usepackage{fourier}
%\usepackage{tgpagella}
%\usepackage{lmodern}
%\usepackage{tgheros}
%\usepackage{newtxtext}

%-----------------------------
%%%%%%  CV STARTS HERE  %%%%%%

\begin{document}




%----------HEADING-----------------

\begin{center}
    \textbf{{\LARGE Harsh Bordekar}} \\
    \textbf{Composite Simulation \& Technology Development Engineer | PhD Candidate at DLR \& TU Braunschweig} \\
    \textbf{CFRP Process \& Thermal Simulation, Manufacturing Quality, Research-to-Production (TRL)}
    
    \vspace{0.5em}
    
    \small
    \begin{tabular}{*{3}{l@{\hspace{2.5em}}}}
        \faIcon[regular]{envelope} \href{mailto:harshbordekar@gmail.com}{harshbordekar@gmail.com} & 
        \faIcon{linkedin} \href{https://www.linkedin.com/in/harsh-s-bordekar/}{Harsh's LinkedIn} & 
%        \faIcon{github} \href{https://github.com/gagankaushikmanyam}{Gagan's Github} \\
        \faIcon{mobile-alt} +49 (0) 17673888310 &
        \faIcon{map-marker-alt} Braunschweig, Germany & 
        %\faIcon{home} Niedersachsen (Residence) & 
        \faIcon{id-card} Permanent Resident - Germany
    \end{tabular}
\end{center}



\vspace{5pt}

\section{\color{blue}Profile}
Composite research and simulation engineer with 7+ years of experience in CFRP materials, high-fidelity FEM (ABAQUS), and thermal and structural simulation of aerospace-relevant structures. Currently a PhD Candidate at DLR \& TU Braunschweig, developing microstructure-informed multiscale models that link manufacturing-induced defects and variability to composite part quality and performance. Experienced in evaluating material systems and effective properties, quantifying manufacturing-induced defects (fibre waviness, internal porosity) from image and CT data, and running large-scale thermal and structural simulations on HPC. Works within collaborative, publicly-funded research projects and technology-development (TRL) settings, interfacing with partner institutes, coordinating student work, and documenting results to aerospace reporting standards and scientific practice. Strong Python-based automation of simulation and post-processing workflows to accelerate iteration and improve reproducibility.

\vspace{0.2cm}
\textbf{Core Competencies:} CFRP \& Composite Materials, Thermal \& Structural Simulation (ABAQUS/ABAQUS CAE), Manufacturing-Induced Defect \& Quality Analysis, Material System \& Technology Evaluation, Non-linear \& Damage Modelling, Multiscale / Homogenization, Research-Project Work-Package Contribution, Technology Development (TRL), Aerospace-Standard Documentation, Python Automation \& Scientific Computing, HPC (SLURM). \textit{Building AFP/process-simulation and CATIA composite-design skills.}



\section{\color{blue}Professional Experience}


\resumeSubheading
    {Research \& Simulation Engineer (CFRP)}{DLR \& TU Braunschweig}
    {Composite modelling, thermal/structural simulation \& research projects}{Feb $2021$ - Present}
    \resumeItemListStart
    \item (PhD) Developed a physics-based, microstructure-informed multiscale framework for CFRP that links manufacturing-induced defects and spatial variability to composite part quality and performance, propagating effects from ply and laminate to structural scale.
\\
\textbf{Tools: Git, Docker, Python, ABAQUS, FEM, HPC/SLURM}

\item Ran large-scale thermal and structural simulations for the Callisto reusable-rocket housing bay, building models to predict temperature fields and structural response, directly relevant to thermal and process simulation of aerospace components.
\\
\textbf{Tools: ABAQUS, SLURM (HPC), Python, Nastran, Patran}

\item Evaluated composite material systems and effective properties via computational homogenization, and built an Agentic AI tool that predicts homogenized properties from microstructure images, enabling fast material screening without full-scale FEM runs.
\\
\textbf{Tools: Python, Physics-based CNN, Hierarchical RAG, image processing, ABAQUS}

\item Contributed to and coordinated work packages within collaborative, publicly-funded research projects, interfacing with partner institutes and students, transferring research outcomes toward application, and documenting results to aerospace reporting standards and scientific practice.
\\
\textbf{Tools: Python, ABAQUS, LaTeX/technical reporting}

\item Instructed 60+ Master's students in FEM theory and ABAQUS-based CFRP modelling, covering material modelling, Continuum Damage Mechanics (CDM), Progressive Damage Analysis (PDA), and Cohesive Zone Modelling (CZM), and supervised theses on composite damage tolerance and fracture mechanics.
\\
\textbf{Tools: ABAQUS, Automation using Python}

\resumeItemListEnd
	

  \vspace{10pt}


\resumeSubheading
    {Research Intern, Manufacturing Quality \& Defect Characterization}{DLR-Braunschweig}
    {CT-based QA of manufactured aerospace parts}{Nov $2019$ - Jan $2021$}
\resumeItemListStart
    \item Built a CT-based defect-detection and quality-assessment framework for additively manufactured aerospace-grade components, automating classification of internal defects (porosity, inclusions) with interpretable confidence outputs to support quality sign-off.
    \\
    \textbf{Tools: Python (scikit-learn / TensorFlow, SHAP / LIME / Grad-CAM), CT image processing}

    \item Coupled measured defect populations with numerical simulation to establish acceptance limits for manufactured parts, supporting technology qualification and manufacturing-process control.
    \\
    \textbf{Tools: FEM (ABAQUS), Python (NumPy/SciPy)}


\resumeItemListEnd

  \vspace{10pt}



	
\resumeSubheading
    {Student Research Assistant}{Leibniz University Hannover}
    {Institute of Static and Dynamic}{Sep $2018$ - Oct $2019$}
    
    \resumeItemListStart
    

    \item Developed an ML framework to detect and quantify fibre undulations (waviness) in unidirectional CFRP laminates, producing spatially resolved layup-quality metrics that link manufacturing-induced defects in the fibre layup to mechanical performance, directly relevant to automated-layup (e.g. AFP) quality assessment.
    \\
    \textbf{Tools: Matlab, Python (scikit-learn / TensorFlow), image processing, NumPy/SciPy}
    
\resumeItemListEnd
\vspace{10pt}

    \resumeSubheading
    {Design Engineer }{Gujarat Industrial Development Corporation}
    {AK. Alloys}{Mar $2017$ - Mar $2018$}
    \resumeItemListStart
   \item Designed sand cast components with manufacturing tolerances and process constraints, translating functional requirements into production-ready geometries compliant with sand casting process limitations, draft angles, parting line constraints, and dimensional quality standards.
       \\
    \textbf{Tools: CREO, SolidWorks, AutoCAD}

\item Managed witness casting and end-to-end quality assurance processes, ensuring dimensional conformance and material integrity of cast components against engineering specifications and customer acceptance criteria.
\\
\textbf{Tools: CREO, SolidWorks, AutoCAD}

\item Planned and optimized sand casting production schedules based on capacity, material availability, and delivery requirements, improving throughput alignment with manufacturing targets.
\textbf{Tools: SolidWorks, AutoCAD, ERP / production planning tools}
\resumeItemListEnd
        
		



%-------------SKILLS SUMMARY------------------
\renewcommand{\arraystretch}{1.3}
\section{\color{blue}Skills Summary}
\begin{tabular}{|p{8cm}p{10.5cm}|}

\hline
\textbf{Composite Materials \& Manufacturing} &
CFRP \(\vert\) Manufacturing-induced defects \& variability \(\vert\) Fibre-waviness / layup-quality metrics \(\vert\) Material system \& technology evaluation \\

\hline
\textbf{Simulation (ABAQUS)} &
ABAQUS / ABAQUS CAE \(\vert\) Thermal \& structural FEM \(\vert\) Non-linear \& damage modelling \(\vert\) Multiscale / homogenization \\

\hline
\textbf{Numerical Simulations} &
Anisotropic material modeling and FEM  \(\vert\) Fracture Mechanics  \(\vert\) Progressive damage analysis \(\vert\) Verification \& Validation \\

\hline
\textbf{ Programming Languages} &
Python  \(\vert\) C++ \(\vert\) FORTRAN \(\vert\) MATLAB \\
\hline

\textbf{ Data Analysis \& Scientific Computing} &
NumPy \(\vert\) Pandas \(\vert\) Matplotlib \(\vert\) Seaborn \(\vert\) Plotly \(\vert\) SALib \(\vert\) SciPy \\
\hline

\textbf{ ML \& AI} &
Scikit-learn \(\vert\) TensorFlow \(\vert\) PyTorch \\
\hline

\textbf{LLMs \& Generative AI} &
Hugging Face Transformers (Fine-tuning) \(\vert\) Prompt Engineering \(\vert\) \\ & Embeddings \(\vert\) RAG \(\vert\) LangChain \(\vert\) LangGraph \(\vert\) MCP \\

\hline

\textbf{ APIs \& Backend Development} &
RESTful APIs \(\vert\) FastAPI \(\vert\) JSON\\
\hline


\textbf{ Version Control \& Reproducibility} &
Git \(\vert\) Docker \\
\hline

\textbf{ Application Development} &
VS Code \(\vert\) JupyterLab \(\vert\) Google Colab \(\vert\) Streamlit \(\vert\) Gradio\\
\hline
\textbf{ Languages} &
English (C1) \(\vert\) German (B1, working toward B2) \\
\hline
\end{tabular}

\vspace{0.2cm}
%\small{\textit{Keywords: Machine Learning, Applied ML, Model Training, Model %Evaluation, Feature Engineering, Inference Pipelines, Reproducible ML}}
\vspace{5pt}

%-----------Awards-----------------

%-----------EDUCATION-----------------
\section{\color{blue}Awards and Education}

\resumeSubheading
  {\textcolor{green!50!black}{\faAward} \textbf{2nd Prize: Best Paper Award}}
  {DLR (German Aerospace Center)}{}{2026}
\resumeSubheading
  {\textcolor{green!50!black}{\faAward} \textbf{3rd Prize: Best Paper Award}}
  {DLR (German Aerospace Center)}{}{2024}
\vspace{4pt}

\resumeSubheading
  {\textcolor{green!50!black}{M.Sc.} Computational Methods in Engineering}
  {Leibniz University Hannover, Germany}{}{April 2018 -- Dec 2020}
  \vspace{2pt}
\textbf{Focus:} Numerical Methods, Scientific Machine Learning, Deep Learning
\vspace{4pt}
\section{\color{blue}Selected Publications}

\begin{enumerate}

  \item \textbf{Bordekar, H.}, et al.
  ``Microstructure-informed Probabilistic Multiscale Modelling of CFRP via Machine-Learning-Derived Stochastic Volume Elements.''
  \textit{Manuscript in preparation}, 2026.

  \item \textbf{Bordekar, H.}, et al.
  ``MicroStructAgent: A Physics-Informed Agentic AI Framework for Computational Homogenization and Quality Assurance of Fiber Composite Laminates.''
  \textit{Manuscript in preparation}, 2026.

  \item \textbf{Bordekar, H.}, et al.
  ``Microstructure Characterization of Fiber Composite Laminate Using eXplainable Artificial Intelligence.''
  \textit{Journal of Composite Materials}, vol. 60, no. 1, 2026, pp. 3--25.
  \href{https://journals.sagepub.com/doi/pdf/10.1177/00219983251345559}{\textcolor{blue}{[Link]}}

  \item \textbf{Bordekar, H.}, et al.
  ``eXplainable Artificial Intelligence for Automatic Defect Detection in Additively Manufactured Parts Using CT Scan Analysis.''
  \textit{Journal of Intelligent Manufacturing}, vol. 36, no. 2, 2025, pp. 957--974.
  \href{https://link.springer.com/article/10.1007/s10845-023-02272-4}{\textcolor{blue}{[Link]}}

\end{enumerate}
\vspace{2pt}

\section{\color{blue}References}

\begin{tabular}{@{} p{6cm} p{10cm}@{}}
\textbf{\href{https://www.dlr.de/de/sy/ueber-uns/abteilungen-dlr-sy/dlr-sy-funktionsleichtbau}{Prof. Dr. Christian Huehne}} & Abteilungsleiter Funktionsleichtbau. Deutsches Zentrum für Luft- und Raumfahrt. Institut für Systemleichtbau \\
\textbf{\href{https://www.tu-braunschweig.de/ima/team/head-of-institute/prof-dr-ing-oliver-voelkerink}{Prof. Dr. Oliver Voelkerink}} & Head of Institute., IMA, TU-Braunschweig\\
\textbf{\href{https://www.linkedin.com/in/nicola-cersullo-41772b131 }{
Dr. Nicola Cersullo}} & R\&T Project Manager @ Airbus Hamburg \\
\end{tabular}

\end{document}

